% !TeX root = ./main.tex
\chapter{广义特征分解}

初学者接触广义特征分解，往往找不到定义的动机。直到现在，我也不知道自己是理解了它还是接受了它（或者还压根没有理解\UseVerb{cry}）

在这部分的学习过程中，我们头顶上难免冒出几朵疑云：
\begin{enumerate}
    \item 为什么一个因式分解块\((\lambda I - T)^{r} \)就恰好是一个不变子空间？
    \item 这些块乘起来构成的式子怎么恰好是行列式\(\det (\lambda I - T)\)？
\end{enumerate}
\section{广义特征向量和幂零算子}

\subsection{小淘气向量}
\begin{theorem}
    设 \(T \in \L(V)\). 那么
    \begin{itemize}
        \item \(\{0\} = \nullspace T^0 \subseteq \nullspace T^1 \subseteq \dots \subseteq \nullspace T^k \subseteq \nullspace T^{k+1} \subseteq \dots\)
        \item 若\(\nullspace T^m = \nullspace T^{m+1}\)，
            则\(\nullspace T^m = \nullspace T^{m+1} = \nullspace T^{m+2} = \nullspace T^{m+3} = \cdots\)
        \item \(\nullspace T ^ {\dim V} = \nullspace T ^ {\dim V + 1} = \nullspace T ^ {\dim V + 2} = \dots\)
    \end{itemize}
\end{theorem}
\begin{proof}

\end{proof}
\begin{theorem}[Fitting引理]
    \[
        V = \nullspace T^{\dim V}\oplus\range T^{\dim V}
    \]
\end{theorem}
\begin{proof}

\end{proof}
\subsection{广义特征向量}
还记得线性映射基本定理里\(V\neq \range T \oplus \nullspace T\)吗？这是因为对于一般的映射\(T\)，
\(\range T\)里总会有小淘气向量\(v=Tu\)，看起来\(v\neq 0\)，
但\(Tv\)却跳到了\(\nullspace T\)里，导致\(Tu\neq 0\)但\(T^{2}u=0\)。
这导致整个空间不能被分成直和。但如果反复应用\(T\)直到零空间收敛，
就能有\(V=\range T^{\dim V}\oplus \nullspace T^{\dim V}\)。

这有什么用呢？回想一下之前我们对算子的分解。我们的梦中算子最好能分成一个个小算子。
即：\[
    V=V_{1}\oplus V_{2}\oplus \dots V_{n}
\]
同时\(V_{i}\)还得是\(T\)的不变子空间。

这实际上是两个要求：不变子空间+直和。

目前，我们做到了：
\begin{itemize}
    \item 在\(\C\)上，任何算子都有上三角形式。也就是层层递增的不变子空间
    \item 对于最小多项式无重根的多项式，还能分解成特征子空间的直和
\end{itemize}
但这离我们分解所有算子的梦想还相差甚远。

我们仔细观察证明过程，不难看出，无重根的证明思路大抵长这样：
\begin{enumerate}
    \item 考虑最小多项式里的一项\(U=\range(T-\lambda_{m}I)\)，受限算子\(T|_{U}\)有特征向量的基。
    \item 在\(\nullspace(T-\lambda_{m}I)\)上\(T\)自然也有特征向量的基
    \item 不同特征值对应基线性无关。
        因此\(V = \range(T-\lambda_{m}I)\oplus\nullspace(T-\lambda_{m}I) \)。
        从而\(T\)有特征向量构成的基，可以对角化。
\end{enumerate}

不难看出，
罪魁祸首就在于一般\(V \neq  \range(T-\lambda_{m}I)\oplus\nullspace(T-\lambda_{m}I)\)，
这就是线性映射基本定理给我们下的绊子！

怎么绕过这个问题呢？还是得借用我们刚刚看到的Fitting引理，
我们把空间分解成\(V=\range T^{\dim V}\oplus \nullspace T^{\dim V}\)就好。

\begin{definition}
    设 \(\lambda\) 是 \(T\) 的特征值. 称向量 \(v \in V\) 是 \(T\) 对应于 \(\lambda\) 的广义特征向量，
    若 \(v \neq 0\) 且对某个正整数 \(k\) 有 \((T - \lambda I)^k v = 0\).
\end{definition}

这样一来，任意复向量空间算子都有广义特征向量的基：
\begin{theorem}
    设 \(\F = \C\) 且 \(T \in \L(V)\). 那么存在由 \(T\) 的广义特征向量构成的 \(V\) 的基.
\end{theorem}

\begin{proof}

\end{proof}

这组基完美符合我们的要求：
\begin{definition}
    \(T\) 对应于 \(\lambda\) 的广义特征空间
    定义为：\(G(\lambda ,T) = \{v\in V:(T - \lambda I)^k v = 0,k\) 为某正整数\}
\end{definition}
不难看出\(G(\lambda,T) = \nullspace T^{\dim V} \)。

\begin{theorem}
    设 \(\F = \C\) 且 \(T \in \L(V)\). 令 \(\lambda_1, \dots, \lambda_m\) 是 \(T\) 的所有互异特征值. 那么
    \begin{itemize}
        \item 对每个 \(k = 1, \dots, m\), \(G(\lambda_k, T)\) 在 \(T\) 下是不变的;
        \item \(V = \bigoplus_ {k = 1}^{m} G (\lambda_{k}, T)\).
    \end{itemize}
\end{theorem}
\begin{proof}

\end{proof}
% todo:证明 p253
% 还有一些无关紧要的东西，比如广义特征向量线性无关等证明，后续补充

% 这时一个子空间上的算子，自然就能分解成幂零部分和不变部分
% 幂零部分的性质什么的，见第一节

% todo: 问题解答
% 1. 简单来说range(p(T))是T的不变子空间，而且任意算子在任意子空间上都有特征值。
% 借此我们把算子分解成null lam I -T（实际上是null lam T - I + range lam T - I，
% 我的值域就是别人的零空间），因此一个块就是一个子空间。至于为什么要研究多项式可见第5章
% 2. 书上的证明是证明代数重数>=几何重数，在上三角分解里主对角线lam次数一定小于子空间维数（实际上就是等于，因为直和）。于是行列式就是这堆东西乘一起。
% 但我觉得这个证明并不本质，可见cherry studio笔记。某种意义上在域论和最小多项式的概念上，算子的广义特征分解就是中国剩余定理

% todo: 后续若当标准型什么的
% 若当标准型的计算
% 还记得大一下学期的高代月考，第一题就是计算矩阵的若当标准型。可惜当时学业不精，第一题就没写出来\UseVerb{sad}，考试后才知道有若当链这么一个东西。
% https://math.stackexchange.com/questions/4495/the-intuition-behind-generalized-eigenvectors
